\documentclass{article}
\usepackage{relsize}
\usepackage{amssymb}
\usepackage{geometry}

% Reduce margins
\geometry{
    a4paper,
    margin=1in % You can adjust this value (e.g., 0.75in, 2cm, etc.)
}

\newcommand{\subtitlerelsize}{1} %relative size: integer value
\newcommand{\subtitlelinesep}{0.2em} %line separation: a LaTeX length

\title{Evaluationsbogen Agenten-Frameworks}
\date{}

\begin{document}
    
\maketitle
    
\vspace{-1.5cm}

\begin{center}
    \textbf{Framework:} y
    \textbf{Testzeitraum:} x Tage
\end{center}

\section*{Teil A}

\newcommand{\evalquestion}[2]{
    #1
    \begin{flushleft}
        \ifnum#2=1 $\boxtimes$ Stimme gar nicht zu \else $\square$ Stimme gar nicht zu \fi \quad
        \ifnum#2=2 $\boxtimes$ Eher nicht \else $\square$ Eher nicht \fi \quad
        \ifnum#2=3 $\boxtimes$ Neutral \else $\square$ Neutral \fi \quad
        \ifnum#2=4 $\boxtimes$ Eher ja \else $\square$ Eher ja \fi \quad
        \ifnum#2=5 $\boxtimes$ Stimme voll zu \else $\square$ Stimme voll zu \fi
    \end{flushleft}
}

\noindent
\evalquestion{1. Ich denke, dass ich dieses Framework häufiger verwenden möchte.}{0}

\noindent
\evalquestion{2. Ich fand das Framework unnötig komplex.}{0}

\noindent
\evalquestion{3. Ich dachte, das Framework war einfach zu bedienen.}{0}

\noindent
\evalquestion{4. Ich denke, dass ich die Unterstützung eines technischen Supports brauche, um dieses Framework nutzen zu können.}{0}

\noindent
\evalquestion{5. Ich fand, die verschiedenen Funktionen in diesem Framework waren gut integriert.}{0}

\noindent
\evalquestion{6. Ich dachte, dass dieses Framework nicht konsistent genug war.}{0}

\noindent
\evalquestion{7. Ich würde mir vorstellen, dass die meisten Informatiker sehr schnell lernen würden, dieses Framework zu benutzen.}{0}

\noindent
\evalquestion{8. Ich fand dieses Framework sehr umständlich zu benutzen.}{0}

\noindent
\evalquestion{9. Ich habe mich sehr selbstsicher gefühlt, dieses Framework zu verwenden.}{0}

\noindent
\evalquestion{10. Ich musste eine Menge lernen, bevor ich mit diesem Framework loslegen konnte.}{0}

\section*{Teil B}

\noindent
\evalquestion{1. Die technische Dokumentation des Frameworks ist vollständig, aktuell und verständlich.}{0}

\noindent
2. Welche wichtigen Informationen waren in der Dokumentation nicht oder nur schwer
auffindbar? \\

\noindent
3. Gab es Aspekte, die sich im praktischen Einsatz als besonders durchdacht oder
problematisch erwiesen haben – unabhängig von der Dokumentation? \\

\noindent
4. Welche Lösungen oder Workarounds mussten implementiert werden, um bestimmte
Anforderungen oder Anwendungsfälle umzusetzen? \\

\noindent
5. Gab es Funktionen oder Eigenschaften des Frameworks, die Sie positiv oder negativ
überrascht haben? \\

\noindent
\evalquestion{6. Wie gut lässt sich das von Ihnen mit dem Framework erstellte System (auch in kleinerem
Maßstab) im Nachhinein ändern oder erweitern?}{0}

\noindent
\evalquestion{7. Wie gut wäre das von Ihnen erstellte System mit vertretbarem Aufwand in einem
produktiven Szenario einsetzbar?}{0}

\noindent
8. Wie haben Sie das System hauptsächlich erstellt?
\begin{flushleft}
    $\square$ Überwiegend durch Programmierung \quad
    $\square$ Über Konfiguration (z. B. YAML, JSON) \quad
    $\square$ Über grafische Benutzeroberfläche (GUI) \quad
\end{flushleft}

\noindent
9. Wie schätzen Sie den Aufwand und die Flexibilität dieser Art der Systemerstellung
( siehe Frage 8 ) ein?



\end{document}